\documentclass[10pt,a4paper]{amsart}
\usepackage[margin=19mm]{geometry}
\usepackage{amsmath,amssymb,mathtools}
\usepackage[scr=rsfs,cal=euler]{mathalfa}
\usepackage{xfrac}
\usepackage[unicode,hidelinks,bookmarks=false]{hyperref}
\newcommand{\ie}{i.e.\ }
\newcommand{\cf}{cf.\ }
\newcommand{\resp}{respectively\ }
\newcommand{\Subsection}{\subsection}
\providecommand{\nobreakdash}{\nobreak\hbox{-}}
\setlength{\emergencystretch}{2em}
\multlinegap0pt
\usepackage{amssymb}
\usepackage{bbm}
\usepackage{xcolor}
\usepackage{bm}
\usepackage[shortlabels]{enumitem}
\usepackage{tikz}
\usepackage{mathtools}
\usepackage{natbib}
\newtheorem{lemma}{Lemma}
\newcommand{\CM}{\H_\mu}
\renewcommand{\H}{\mathcal{H}}
\newcommand{\R}{\mathbb{R}}
\renewcommand{\ge}{\geqslant}
\renewcommand{\le}{\leqslant}
\newcommand{\nn}{\nonumber}
\title{Uniqueness for stochastic differential equations in Hilbert spaces with irregular drift}
\author{Lukas Anzeletti, Oleg Butkovsky, M\'at\'e Gerencs\'er, Alexander Shaposhnikov}
\date{Source: arXiv:2512.25003v3 (CC BY 4.0)}
\begin{document}
\maketitle

\noindent\emph{Source and licence.} Uniqueness for stochastic differential equations in Hilbert spaces with irregular drift. Version arXiv:2512.25003v3 (CC BY 4.0), distributed by arXiv under the Creative Commons Attribution 4.0 International licence, http://creativecommons.org/licenses/by/4.0/, as recorded in the arXiv licence metadata for this exact version and shown on the abstract page https://arxiv.org/abs/2512.25003v3. Full licence notice: \texttt{licenses/arxiv-2512.25003-CC-BY-4.0.txt}.\par\smallskip

\noindent\textit{Editorial scope.} Counted unit: the article's lemma lem:averagingproperties (source0.tex lines 552--559). The article's complete original proof of it is delivered with it (source0.tex lines 560--630, one \texttt{proof} environment of 71 lines). No mathematics is written, repaired or re-derived: the statement and the proof are the article's own lines, in the article's own notation, and no retained line is substituted or re-flowed.  Retained uncounted as the source-local context the proof needs: lines 547--549.  Nothing was omitted from any retained span, so the package carries no omission record.  No retained line contains a citation, so the wrapper carries no bibliography and no bibliography entry.  No retained line contains a figure environment, an image-inclusion command or drawing code, so no image is delivered and the package declares no asset dependency.  Editorial formatting only, in addition: an AMS \texttt{amsart} preamble that restates the article's own declarations and macros used by the retained lines, namely the environments \texttt{lemma} on separate counters, together with the standard packages those lines need.  The retained environments print in this document as Lemma 1 in order of appearance, whereas the article prints the same items with the numbering of its own full document.  The complete native source of the article as distributed in the arXiv e-print for this exact version is bundled unchanged as \texttt{sources/191954/source0.tex}.
\begin{equation}\label{eq:R-N}
\frac{d T_h^\# \mu}{d \mu}(x)=\exp\Big\{\langle h^*,x\rangle_H-\frac{1}{2}\|h\|_{\CM}^2\Big\},\quad x\in H.
\end{equation}

\begin{lemma} \label{lem:averagingproperties}
There exists a constant $C>0$ such that for any bounded measurable function $f\colon H\to H$ and any $h_1,h_2,h_3 \in \mathcal{H}_{\mu}$, one has
\begin{align}
&\Big\|\int_H f(x+h_1)-f(x+h_2) \mu(dx)\Big\|_{H}\leqslant C\|f\|_\infty \|h_1-h_2\|_{\CM} \label{eq:2point},\\
&\Big\|\int_H f(x+h_1)-f(x+h_2)-f(x+h_3)+f(x+h_3+h_2-h_1) \mu(dx)\Big\|_{H} \nonumber\\
    &\qquad\le C \|f\|_\infty \|h_3-h_1\|_{\CM} \|h_2-h_1\|_{\CM}. \label{eq:4point}
\end{align}
\end{lemma}

\begin{proof}
First, let us prove  \eqref{eq:2point}.
It suffices to consider the case $h_2=0$; the case $h_2\neq 0$ then easily follows by shifting $f$. Note that if $\|h_1\|_{\CM}\geqslant 1$, then \eqref{eq:2point} holds trivially. Henceforth we assume $\|h_1\|_{\CM}\le1$. Denote $\zeta_h(x):=\langle h^*,x\rangle_H-\frac{1}{2}\|h\|_{\CM}^2$. 
Using \eqref{eq:R-N} and the inequality
\begin{equation}\label{eq:very-triv}
     |e^{a}-1|\le (1\vee e^{a})|a|
\end{equation}
valid for any $a\in \R$, we derive
\begin{align}\label{2pointmany}
\left\|\int_H f(x+h_1)-f(x)\mu(dx)\right\|_{H}&=\left\|\int_H f(x)  T_{h_1}^\#\mu (dx)-\int_H f(x)\mu(dx)\right\|_{H}\nn\\
&=\left\|\int_H f(x)\Bigl(\frac{d T_{h_1}^\# \mu}{d \mu}(x)-1\Bigr) \mu(dx)\right\|_{H}\nn\\
&\leqslant \|f\|_\infty \int_H (1\vee e^{\zeta_{h_1}(x)}) |\zeta_{h_1}(x)|\mu(dx)\nn\\
&\leqslant \|f\|_\infty \|1\vee e^{\zeta_{h_1}(\cdot)}\|_{L^2(H,\mu)} \|\zeta_{h_1}(\cdot)\|_{L^2(H,\mu)}.
\end{align}


Recalling the variance of $\zeta_{h_1}$, we see that 
\begin{equation}\label{eq:constant1}
\|1\vee e^{\zeta_{h_1}(\cdot)}\|_{L^2(H,\mu)}^2
\le 1+\exp(\|h_1\|^2_{\CM})\le 4,
\end{equation}
since we consider the case $\|h_1\|_{\CM}\le1$. 
Again using the same arguments and the inequality $(a+b)^2\le 2 a^2 +2 b^2$ valid for any $a,b\in\R$ we get,
\begin{equation}\label{eq:constant2}
\|\zeta_{h_1}(\cdot)\|_{L^2(H,\mu)}^2\le 2 
 \|\langle h_1^*,\cdot\rangle \|_{L^2(H,\mu)}^2 +\frac12\|h_1\|_{\CM}^4 \le 3\|h_1\|_{\CM}^2.
\end{equation}
Combining this with \eqref{eq:constant1} and substituting into \eqref{2pointmany}, we get  \eqref{eq:2point}.

Now we move on to the proof of  \eqref{eq:4point}. We note again that it suffices to prove this inequality  for $h_1=0$ as the case $h_1\neq 0$ then easily follows by shifting $f$. If $\|h_2\|_{\CM}\ge1$ and $|h_3\|_{\CM}\ge 1$, \eqref{eq:4point} is immediate.
It remains to treat the cases $\|h_3\|_{\CM}\leqslant 1 \leqslant \|h_2\|_{\CM}$ (the case of both inequalities reversed follows by symmetry) and $\|h_3\|_{\CM}\le 1$, $\|h_2\|_{\CM} \leqslant 1$.

First, consider the case $\|h_3\|_{\CM}\leqslant 1 \leqslant \|h_2\|_{\CM}$.
Using \eqref{eq:2point} for the functions $f$ and $T_{h_2}f$, we get
\begin{align*}
    \Big\|\int_H &f(x)-f(x+h_2)-f(x+h_3)+f(x+h_3+h_2) \mu(dx)\Big\|_{H} \\ &=\Big\|\int_H f(x)-f(x+h_3)-(T_{h_2}f(x)-T_{h_2}f(x+h_3)) \mu(dx) \Big\|_{H}\\
    &\leqslant C(\|f\|_\infty + \|T_{h_2}f\|_\infty)\|h_3\|_{\CM}\le C \|f\|_\infty \|h_2\|_{\CM}\|h_3\|_{\CM},
\end{align*}
with a universal constant $C>0$. This yields \eqref{eq:4point}.

Next, consider the case $\|h_2\|_{\CM}\le1$, $\|h_3\|_{\CM} \leqslant 1$. 
We have
\begin{align}\label{bigstuff}
&\Big\|\int_H f(x)-f(x+h_2)-f(x+h_3)+f(x+h_2+h_3) \mu(dx)\Big\|_{H}\nn\\
&\qquad=\Big\|\int_H f(x)\mu(dx) -\int_H f(x)T_{h_2}^\#\mu(dx) -\int_H f(x)T_{h_3}^\#\mu(dx) +\int_H f(x)T_{h_2+h_3}^\#\mu(dx) \Big\|_{H}\nn\\
&\qquad= \Big\|\int_H f(x)(1-e^{\zeta_{h_2}(x)}-e^{\zeta_{h_3}(x)}+e^{\zeta_{h_2+h_3}(x)})\mu(dx)\Big\|_{H}\nn\\
&\qquad\leqslant \Big\|\int_H f(x)(e^{\zeta_{h_2}(x)}-1)(e^{\zeta_{h_3}(x)}-1)\mu(dx)\Big\|_{H}\nn\\
&\qquad{\phantom{\le}} +\Big\|\int_H f(x)\Big(e^{\zeta_{h_2+h_3}(x)}-e^{\zeta_{h_2}(x)}e^{\zeta_{h_3}(x)}\Big) \mu(dx)\Big\|_{H}=:I_1+I_2.
\end{align}
The first summand $I_1$ is treated similarly to the proof of \eqref{eq:2point}. Using \eqref{eq:very-triv} and H\"older's inequality, we get 
\begin{align*}
I_1&\le \|f\|_\infty\Big\|\int_H  (1\vee e^{\zeta_{h_2}(x)}) |\zeta_{h_2}(x)|(1\vee e^{\zeta_{h_3}(x)}) |\zeta_{h_3}(x)| \mu(dx)\Big\|_{H}\\
&\le\|f\|_\infty \prod_{i=2,3} \|1\vee e^{\zeta_{h_i}(\cdot)}\|_{L^4(H,\mu)} \|\zeta_{h_i}(\cdot)\|_{L^4(H,\mu)}.
\end{align*}
Similarly to \eqref{eq:constant1} and \eqref{eq:constant2}, we derive for $i=2,3$
\begin{equation*}
\|1\vee e^{\zeta_{h_i}(\cdot)}\|_{L^4(H,\mu)}^4\le C;\qquad
\|\zeta_{h_i}(\cdot)\|_{L^4(H,\mu)}^4\le C \|h_i\|_{\CM}^4   
\end{equation*}
for a universal constant $C>0$. We can conclude
\begin{equation}\label{ivangoal}
    I_1\leqslant C\|f\|_\infty\|h_2\|_{\CM}\|h_3\|_{\CM}. 
\end{equation}
As for $I_2$, we have $e^{\zeta_{h_2+h_3}(x)}-e^{\zeta_{h_2}(x)}e^{\zeta_{h_3}(x)}=e^{\zeta_{h_2+h_3}(x)}(1-e^{\langle h_2,h_3\rangle_{\CM}})$. Therefore
\begin{align*}
I_2&=|1-e^{\langle h_2,h_3\rangle_{\CM}}|\,\Big\|\int_H f(x)T_{h_2+h_3}^\# \mu(dx)\Big\|_{H}\\
&\le  e\|f\|_\infty|\langle h_2,h_3\rangle_{\CM}|\leqslant e \|f\|_\infty \|h_2\|_{\CM} \|h_3\|_{\CM}.
\end{align*}
Combining this with \eqref{ivangoal} and \eqref{bigstuff}, we get 
\eqref{eq:4point}.
\end{proof}

\end{document}
